\begin{workedexample}[Worked Example: Gradient slew rate]
A gradient must ramp quickly to keep echo times short. Slew rate is
$\text{SR} = G_{max}/t_{rise}$. Ramping to $G_{max} = 40\ \text{mT/m}$ in
$t_{rise} = 200\ \mu\text{s}$ gives
\[ \text{SR} = \frac{0.04}{200\times10^{-6}} = 200\ \text{T/m/s}. \]
Faster ramps shorten sequences but are capped by peripheral nerve stimulation
(PNS), which responds to $dB/dt$.
\end{workedexample}

\begin{keytakeaways}
\begin{itemize}[leftmargin=1.4em,itemsep=2pt]
\item Gradients make resonance position-dependent: $f(z) = \gamma(B_0 + Gz)$.
\item Slew rate $= G_{max}/t_{rise}$; higher slew shortens TE and TR.
\item PNS (from $dB/dt$) and eddy currents cap how fast gradients can switch.
\end{itemize}
\end{keytakeaways}

\begin{exercises}
\item Slew rate to reach $30\ \text{mT/m}$ in $150\ \mu\text{s}$?
\item What physiological effect limits maximum slew rate?
\item Faster gradient switching allows shorter what?
\end{exercises}

\furtherreading{Nishimura~\cite{nishimura} (ch.~5); Haacke~\cite{haacke}.}
